\documentclass[12pt, a4paper]{article}
\usepackage{iftex}
\ifPDFTeX
  \usepackage{cmap}
  \usepackage[T2A]{fontenc}
  \usepackage[utf8]{inputenc}
  \usepackage{tempora}
\else
  \usepackage{fontspec}
  \IfFontExistsTF{Times New Roman}
    {\setmainfont{Times New Roman}}
    {\setmainfont{CMU Serif}}
  \IfFontExistsTF{Consolas}
    {\setmonofont{Consolas}[Scale=MatchLowercase]}
    {\IfFontExistsTF{Courier New}
      {\setmonofont{Courier New}}
      {\setmonofont{CMU Typewriter Text}}}
\fi
\usepackage[english,russian]{babel}

\usepackage[a4paper, margin=2.54cm]{geometry}
\setlength{\parindent}{0pt}
\setlength{\parskip}{4pt plus 1pt}
\frenchspacing
\setlength{\emergencystretch}{3em}
\raggedbottom

\usepackage{titlesec}
\setcounter{secnumdepth}{0}
\titleformat{\section}{\large\bfseries}{}{0pt}{}
\titleformat{\subsection}{\normalsize\bfseries}{}{0pt}{}
\titlespacing*{\section}{0pt}{14pt plus 4pt minus 2pt}{4pt plus 2pt}
\titlespacing*{\subsection}{0pt}{10pt plus 2pt minus 2pt}{3pt plus 1pt}

\usepackage[table]{xcolor}
\usepackage{graphicx}
\usepackage{tabularx}
\usepackage{float}
\usepackage{caption}
\DeclareCaptionLabelSeparator{spacedemdash}{~--- }
\captionsetup{labelsep=spacedemdash, font=small, singlelinecheck=false, justification=raggedright}
\captionsetup[table]{position=top, skip=4pt}
\DeclareCaptionType[fileext=lol]{listing}[Листинг][Листинги]
\captionsetup[listing]{position=top, skip=4pt, hypcap=false}
\addto\captionsrussian{\renewcommand{\figurename}{Рисунок}\renewcommand{\tablename}{Таблица}}
\newcommand{\thead}[1]{\multicolumn{1}{|c|}{\cellcolor{gray!20}\textbf{#1}}}
\newcolumntype{L}{>{\raggedright\arraybackslash}X}
\newcolumntype{P}[1]{>{\raggedright\arraybackslash}p{#1}}

\usepackage{fvextra}
\usepackage{needspace}
\fvset{breaklines=true, breakanywhere=true,
       frame=leftline, framerule=1.5pt, framesep=6pt,
       rulecolor=\color{gray!60}, xleftmargin=4pt}
\DefineVerbatimEnvironment{code}{Verbatim}{fontsize=\small}
\DefineVerbatimEnvironment{terminal}{Verbatim}{fontsize=\footnotesize}
\newcommand{\inputterminal}[2][]{\VerbatimInput[fontsize=\footnotesize,#1]{#2}}
\ifPDFTeX
  \usepackage{listingsutf8}
  \lstset{inputencoding=utf8/cp1251, basicstyle=\fontencoding{T2A}\selectfont\ttfamily\small,
          breaklines=true, breakatwhitespace=false, columns=fullflexible,
          keepspaces=true, showstringspaces=false, frame=leftline,
          framerule=1.5pt, framesep=6pt, xleftmargin=4pt,
          rulecolor=\color{gray!60}}
  \newcommand{\inputcode}[2][]{\lstinputlisting[#1]{#2}}
  \newcommand{\inputappendixcode}[2][]{\lstinputlisting[basicstyle=\linespread{0.90}\fontencoding{T2A}\selectfont\ttfamily\footnotesize,#1]{#2}}
\else
  \newcommand{\inputcode}[2][]{\VerbatimInput[fontsize=\small,#1]{#2}}
  \newcommand{\inputappendixcode}[2][]{\VerbatimInput[fontsize=\footnotesize,baselinestretch=0.90,#1]{#2}}
\fi
\newenvironment{listingblock}[1]
  {\Needspace{8\baselineskip}\par\medskip\captionof{listing}{#1}\nopagebreak\vspace{-\topsep}\vspace{-\parskip}}
  {\par}

\newcounter{appendixno}
\newcommand{\appendixsection}[2]{\clearpage\section{Приложение #1. #2}\stepcounter{appendixno}\setcounter{listing}{0}\renewcommand{\thelisting}{#1.\arabic{listing}}\renewcommand{\theHlisting}{app\arabic{appendixno}.\arabic{listing}}}

\usepackage[hidelinks]{hyperref}
\usepackage{xurl}

\newcommand{\hint}[1]{\leavevmode{\color{gray}\itshape #1}}
\newcommand{\placeholder}[2][4.5cm]{\fbox{\parbox[c][#1][c]{0.95\linewidth}{\centering\hint{#2}}}}


\newcommand{\ReportLab}{4}
\newcommand{\ReportStudent}{Новак Дмитрий Евгеньевич}
\newcommand{\ReportGroup}{М8О-201БВ-25}
\newcommand{\ReportTeacher}{Миронов Евгений Сергеевич}
\newcommand{\ReportDate}{12.11.2026}
\newcommand{\ReportYear}{2026}
\newcommand{\ReportCourse}{Операционные системы и архитектура компьютера}

\begin{document}
% Титульный лист. Данные берутся из команд \Report... в report.tex.
\begin{titlepage}
  \centering\large

  МОСКОВСКИЙ АВИАЦИОННЫЙ ИНСТИТУТ\\
  (НАЦИОНАЛЬНЫЙ ИССЛЕДОВАТЕЛЬСКИЙ УНИВЕРСИТЕТ)\par
  \vspace{18pt}
  Институт №8 «Компьютерные науки и прикладная математика»\\
  Кафедра 806 «Вычислительная математика и программирование»\par

  \vspace{\stretch{1.2}}

  {\bfseries Лабораторная работа №\,\ReportLab{} по курсу\\
  «\ReportCourse»\par}

  \vspace{\stretch{1.6}}

  \begin{flushright}
    \begin{tabular}{r@{\hspace{0.6em}}l}
      Выполнил:      & \ReportStudent \\[2pt]
      Группа:        & \ReportGroup   \\[2pt]
      Преподаватель: & \ReportTeacher \\[2pt]
      Дата:          & \ReportDate    \\[2pt]
      Оценка:        & \rule[-2pt]{4.2cm}{0.4pt} \\
    \end{tabular}
  \end{flushright}

  \vspace{\stretch{1}}

  Москва, \ReportYear
\end{titlepage}
\setcounter{page}{2}


\section{Условие}

\textbf{Цель работы:} освоить принципы работы с файловыми системами и получить практические навыки обмена данными между процессами посредством отображаемых в память файлов (memory-mapped files).

\textbf{Задание:} составить и отладить программу на языке C, работающую с процессами в ОС семейства UNIX. Основной процесс должен создать дочерний процесс для решения задачи. Межпроцессное взаимодействие организуется через отображаемые файлы и системные сигналы. Необходимо обрабатывать системные ошибки, которые могут возникнуть при выполнении программы. Вариант задачи сохраняется из лабораторной работы №2.

\textbf{Вариант 9:} пользователь вводит имя файла в родительский процесс. Родитель открывает файл на чтение и перенаправляет его в стандартный ввод ребёнка. Каждая строка файла содержит произвольное количество чисел типа \texttt{float}. Дочерний процесс последовательно делит первое число на все следующие числа строки. Результат передаётся родителю, который выводит его в свой \texttt{stdout}. При делении на ноль проверка выполняется в дочернем процессе, после чего оба процесса завершают работу.

\section{Метод решения}

Родительская и дочерняя части реализованы отдельными программами: \texttt{lab4\_parent} и \texttt{lab4\_child}. В отличие от лабораторной работы №2, результаты передаются через общую область памяти, связанную с временным файлом. Сигналы используются для уведомления о готовности результата и подтверждения его получения.

Родитель открывает входной файл и создаёт временный файл функцией \texttt{mkstemp()}. Вызов \texttt{ftruncate()} задаёт размер временного файла, равный \texttt{sizeof(float)}. Затем \texttt{mmap()} с флагом \texttt{MAP\_SHARED} создаёт отображение файла в память. В используемой системе размер \texttt{float} равен четырём байтам, поэтому в общей области хранится один результат.

Имя временного файла удаляется через \texttt{unlink()} сразу после его создания. Это не уничтожает используемый файл: он остаётся доступен через открытый дескриптор и отображения. После освобождения всех ссылок ОС удаляет его содержимое, поэтому отдельный файл с результатами в каталоге не остаётся.

После \texttt{fork()} дочерняя ветвь перенаправляет входной файл в \texttt{stdin} через \texttt{dup2()} и выполняет \texttt{execl()}. Дескриптор временного файла наследуется ребёнком и передаётся новой программе как аргумент командной строки. После \texttt{exec} ребёнок создаёт собственное отображение этого же файла. Виртуальные адреса отображений у процессов могут различаться.

Дочерний процесс читает строки через \texttt{getline()}, преобразует числа функцией \texttt{strtof()} и выполняет деление. Пустые строки пропускаются; строка с одним числом возвращает это число. Нулевой делитель проверяется перед операцией деления, а некорректные числа и выход результата за диапазон \texttt{float} приводят к сообщению об ошибке и завершению.

Обмен результатами выполняется в следующем порядке:
\begin{enumerate}
\raggedright
  \item Ребёнок записывает вычисленный результат в общую область: \texttt{*shared = result}.
  \item Ребёнок отправляет родителю \texttt{SIGUSR1} и ожидает подтверждение через \texttt{sigwait()}.
  \item Родитель получает уведомление, читает \texttt{*shared}, выводит значение и отправляет ребёнку \texttt{SIGUSR2}.
  \item После подтверждения ребёнок переходит к следующей строке и может перезаписать общую ячейку.
\end{enumerate}

Такое подтверждение не позволяет ребёнку перезаписать результат до того, как его прочитает родитель. Сами числа не передаются внутри сигналов: сигналы несут только уведомления.

Родитель блокирует \texttt{SIGUSR1} и \texttt{SIGCHLD} до создания ребёнка. Обработчик сигналов устанавливает флаги \texttt{result\_ready} и \texttt{child\_finished}. Основной цикл родителя проверяет эти флаги, а \texttt{sigsuspend()} временно разрешает доставку сигналов и приостанавливает выполнение. Ребёнок блокирует \texttt{SIGUSR2} до первой записи и принимает подтверждения через \texttt{sigwait()}. Ранний сигнал остаётся ожидающим, поэтому уведомление не теряется между проверкой состояния и ожиданием.

При окончании файла ребёнок освобождает отображение и завершается. Родитель получает \texttt{SIGCHLD}, вызывает \texttt{waitpid()} и возвращает код завершения ребёнка. Затем освобождается отображение родителя. При ошибке родителя после запуска ребёнка выполняется принудительное завершение ребёнка и ожидание его статуса.

\section{Описание программы}

\subsection{Структура программной части}

\begin{table}[H]
  \caption{Файлы программной части}
  \begin{tabularx}{\linewidth}{|P{4.2cm}|L|}
    \hline
    \thead{Файл} & \thead{Назначение} \\ \hline
    \path{src/parent.c} & Ввод имени файла, создание отображаемого файла и ребёнка, приём уведомлений, вывод результатов и ожидание завершения. \\ \hline
    \path{src/child.c} & Чтение строк из стандартного ввода, деление чисел, запись результата в общую память и ожидание подтверждения. \\ \hline
    \texttt{CMakeLists.txt} & Сборка исполняемых файлов \texttt{lab4\_parent} и \texttt{lab4\_child} в стандарте C11. \\ \hline
    \path{data/*.txt} & Примеры корректного ввода, отрицательных чисел и деления на ноль. \\ \hline
    \path{report/trace/*.txt} & Полная и сокращённая трассировки выполнения на Ubuntu ARM64. \\ \hline
  \end{tabularx}
\end{table}

\subsection{Основные переменные и функции}

Вся основная логика каждой программы находится в её \texttt{main()}. В родительской программе отдельно определён обработчик \texttt{handle\_signal()}, который устанавливает флаг готовности результата при \texttt{SIGUSR1} и флаг завершения ребёнка при \texttt{SIGCHLD}. Вывод и работа с общей памятью выполняются основным кодом, а не внутри обработчика.

Переменная \texttt{input\_fd} хранит дескриптор входного файла, \texttt{shared\_fd} --- дескриптор отображаемого файла, \texttt{shared} --- адрес отображения. После создания отображения дескриптор можно закрыть: отображение остаётся действительным до \texttt{munmap()}.

Дочерний процесс хранит текущую строку в динамическом буфере \texttt{line}, который переиспользуется функцией \texttt{getline()} и освобождается через \texttt{free()}. Переменная \texttt{result} имеет тип \texttt{float}; \texttt{line\_number} используется в сообщениях об ошибках.

Коды завершения: 0 --- успех, 1 --- общая или системная ошибка, 2 --- деление на ноль, 3 --- некорректные входные числа или выход результата за диапазон \texttt{float}, 4 --- ошибка запуска дочерней программы через \texttt{exec}.

\subsection{Системные вызовы и POSIX-интерфейсы}

\begin{table}[H]
  \caption{Используемые системные вызовы и POSIX-интерфейсы}
  \begin{tabularx}{\linewidth}{|P{5.2cm}|L|}
    \hline
    \thead{Вызов} & \thead{Назначение в программе} \\ \hline
    \texttt{open()} & Открытие входного файла только для чтения. \\ \hline
    \texttt{mkstemp()}, \texttt{unlink()} & Создание временного файла и удаление его имени из каталога. \\ \hline
    \texttt{ftruncate()} & Установка размера отображаемого файла. \\ \hline
    \texttt{mmap()}, \texttt{munmap()} & Создание и освобождение отображения файла в память. \\ \hline
    \texttt{fork()}, \texttt{execl()} & Создание ребёнка и запуск дочерней программы. \\ \hline
    \texttt{dup2()}, \texttt{close()} & Перенаправление входного файла в \texttt{stdin} и закрытие дескрипторов. \\ \hline
    \texttt{sigaction()} & Установка обработчиков \texttt{SIGUSR1} и \texttt{SIGCHLD}. \\ \hline
    \texttt{sigprocmask()} & Блокировка сигналов до их обработки или ожидания. \\ \hline
    \texttt{sigsuspend()}, \texttt{sigwait()} & Ожидание уведомления у родителя и подтверждения у ребёнка. \\ \hline
    \texttt{kill()} & Отправка сигналов другому процессу. \\ \hline
    \texttt{waitpid()} & Получение статуса завершения дочернего процесса. \\ \hline
  \end{tabularx}
\end{table}

\texttt{getline()}, \texttt{strtof()}, \texttt{printf()} и \texttt{fflush()} относятся к стандартной библиотеке C. Имена функций в исходном коде могут отличаться от системных вызовов, которые видны в \texttt{strace}.

\subsection{Ключевые фрагменты кода}

Родитель передаёт дескриптор отображаемого файла строковым аргументом и отображает область размером в один \texttt{float}:

\begin{listingblock}{Подготовка дескриптора и создание отображения (parent.c)}
  \inputcode[firstline=137,lastline=140]{../src/parent.c}
\end{listingblock}

Ребёнок записывает число в общую память, отправляет уведомление и ждёт подтверждения. Операция записи в \texttt{*shared} является обычным обращением к памяти и не требует отдельного вызова \texttt{write()}.

\begin{listingblock}{Передача результата и ожидание подтверждения (child.c)}
  \inputcode[firstline=120,lastline=135]{../src/child.c}
\end{listingblock}

Родитель ожидает сигнал, выводит доступное значение и подтверждает его получение:

\begin{listingblock}{Ожидание и чтение результата (parent.c)}
  \inputcode[firstline=205,lastline=234]{../src/parent.c}
\end{listingblock}

\section{Инструкция по запуску}

Сборка из корня репозитория с помощью CMake:

\begin{terminal}
$ cmake -S lab_04 -B build/lab_04
$ cmake --build build/lab_04
\end{terminal}

Для запуска нужно перейти в каталог сборки: дочерняя программа ищется по относительному пути \texttt{./lab4\_child}.

\begin{terminal}
$ cd build/lab_04
$ ./lab4_parent
Введите путь к файлу: ../../lab_04/data/ok.txt
0.5
20
10
100000
500
0.05
0
\end{terminal}

Первая строка файла содержит \texttt{32 2 2 2 2 2 2}. Результат вычисления:
\[
32 / 2 / 2 / 2 / 2 / 2 / 2 = 0{,}5.
\]
Строка \texttt{10} возвращает число без изменения. При вычислениях используется тип \texttt{float}, а для вывода --- формат \texttt{\%g}.

Для файла с отрицательными числами получен следующий результат:

\begin{terminal}
$ ./lab4_parent
Введите путь к файлу: ../../lab_04/data/negatives.txt
0.1
-100
-1000
\end{terminal}

При нулевом делителе дочерний процесс сообщает об ошибке, а оба процесса завершаются с кодом 2:

\begin{terminal}
$ ./lab4_parent
Введите путь к файлу: ../../lab_04/data/zero.txt
Ошибка в строке 1: деление на ноль.
$ echo $?
2
\end{terminal}

Родитель читает из своего \texttt{stdin} именно имя файла. Чтобы выполнить запуск без ручного ввода пути, можно передать строку через оболочку:

\begin{terminal}
$ printf '../../lab_04/data/ok.txt\n' | ./lab4_parent
\end{terminal}

В данном случае канал создаёт оболочка только для передачи имени файла родителю. Обмен результатами между родителем и ребёнком по-прежнему выполняется через \texttt{mmap}.

\section{Демонстрация работы strace}

Полная трассировка получена в Ubuntu ARM64. Команды выполняются из корня репозитория. Ключ \texttt{-f} включает наблюдение за дочерним процессом, а \texttt{-o} сохраняет системные вызовы и сигналы в файл.

\begin{terminal}
$ mkdir -p lab_04/report/trace
$ cd build/lab_04
$ printf '../../lab_04/data/ok.txt\n' | strace -f -o ../../lab_04/report/trace/strace.txt ./lab4_parent
\end{terminal}

В данной трассировке PID родителя равен 5028, PID ребёнка --- 5029. Ниже приведены выбранные строки из того же запуска; полная трассировка находится в приложении Б.

\begin{listingblock}{Ключевые системные вызовы из strace}
  \inputterminal[baselinestretch=0.95]{trace/strace_short.txt}
\end{listingblock}

Входному файлу соответствует дескриптор 3, временному --- 4. Размер временного файла составляет четыре байта. Флаг \texttt{MAP\_SHARED} обеспечивает общий доступ к содержимому одного файла, хотя адреса отображений различаются. Закрытие дескриптора после \texttt{mmap()} не отменяет отображение.

В этой системе функциям \texttt{fork()}, \texttt{dup2()}, \texttt{execl()} и \texttt{waitpid()} соответствуют системные вызовы \texttt{clone()}, \texttt{dup3()}, \texttt{execve()} и \texttt{wait4()}. Запуск \texttt{lab4\_child} через \texttt{exec} не меняет PID 5029.

Ребёнок записывает первый результат в общую память и отправляет родителю \texttt{SIGUSR1}. Родитель выходит из \texttt{rt\_sigsuspend()}, выводит \texttt{0.5} через \texttt{write(1, ...)} и подтверждает получение сигналом \texttt{SIGUSR2}. Ожиданию \texttt{sigwait()} у ребёнка соответствует \texttt{rt\_sigtimedwait()}.

Вычисление и запись в \texttt{*shared} не являются системными вызовами и отдельно в трассировке не видны. Большая часть других вызовов \texttt{mmap} относится к загрузке библиотек и служебной памяти.

Пометки \texttt{unfinished} и \texttt{resumed} разделяют начало и продолжение одного вызова. Возврат \texttt{EINTR} при доставке сигнала здесь означает нормальное прерывание ожидания.

В конце ребёнок освобождает отображение и завершается. Родитель получает \texttt{SIGCHLD}, забирает статус через \texttt{wait4()} и освобождает своё отображение. Оба процесса завершаются успешно: в трассировке присутствуют две строки \texttt{exited with 0}.

\Needspace{8\baselineskip}
\section{Выводы}

Таким образом, данная лабораторная работа никак кардинально не изменила мою жизнь: я сделал свою версию 3-way handshake(у меня бюджетная 2-way) как в TCP, поработал с mmap, который занимает половину вывода strace, и научил два процесса по очереди пользоваться четырьмя байтами памяти. Также выяснил, что kill не обязательно кого-то убивает, иногда он просто сообщает, что число готово.

\appendixsection{А}{Исходный код программы}

\begin{listingblock}{Файл parent.c}
  \inputappendixcode{../src/parent.c}
\end{listingblock}

\begin{listingblock}{Файл child.c}
  \inputappendixcode{../src/child.c}
\end{listingblock}

\begin{listingblock}{Файл CMakeLists.txt}
  \inputappendixcode{../CMakeLists.txt}
\end{listingblock}

\appendixsection{Б}{Полный вывод strace}

Полный вывод получен при запуске \texttt{lab4\_parent} в Ubuntu ARM64 с параметром \texttt{-f}. Трассировка включает родительский и дочерний процессы, загрузку библиотек, обмен сигналами и завершение.

\begin{listingblock}{Полный вывод strace -f ./lab4\_parent}
  \inputterminal[baselinestretch=0.95]{trace/strace.txt}
\end{listingblock}

\end{document}
